\subsection{借助一族伪度量定义一致结构}

我们在第六章 § 2 第 3 号中已经看到，若对每个实数 \(a > 0\) 用 \(\mathbf{U}_a\) 表示由所有点对 \((x, y)\) （由 \(\mathbf{R}^n\) 的点组成，其相互间的欧氏距离 \(\leqslant a\) ）组成的集，则诸 \(\mathbf{U}_a\) 构成 \(\mathbf{R}^n\) 的一致结构的一个围域基本系（当 \(a\) 取遍 \(> 0\) 的实数集时）。

更一般地，设 \(f\) 是集合 \(\mathbf{X}\) 上的一个伪度量；对每个 \(a > 0\) ，用 \(\mathbf{U}_a\) 表示 \(\overline{f}^{1}([o, a])\) ，我们来证明：当 \(a\) 取遍一切 \(> 0\) 的实数集时，诸 \(\mathbf{U}_a\) 构成 \(\mathbf{X}\) 上某个一致结构的一个围域基本系。公理 \((\mathbf{U}_{\mathbf{I}}^{\prime})\) 由 \((\mathrm{EC}_{\mathbf{I}})\) 得以满足；若 \(a \leqslant b\) ，则 \(\mathbf{U}_a \subset \mathbf{U}_b\) ，因而诸 \(\mathbf{U}_a\) 满足 \((\mathbf{B}_{\mathbf{I}})\) ；由 \((\mathrm{EC}_{\mathbf{II}})\) 有 \(\overline{\mathbf{U}}_a = \mathbf{U}_a\) ，因而 \((\mathbf{U}_{\mathbf{II}}^{\prime})\) 满足；最后，由 \((\mathrm{EC}_{\mathbf{III}})\) 有 \(\overline{\mathbf{U}}_a \subset \mathbf{U}_{2a}\) ，故 \((\mathbf{U}_{\mathbf{III}}^{\prime})\) 满足。{[}参看第二章，§ 1，第 1 号，定义 2{]}。于是我们可以给出下面的定义：

定义 2. 给定伪度量 \(f\) （在集合 \(\mathbf{X}\) 上），由 \(f\) 定义的一致结构是指 \(\mathbf{X}\) 上以诸集 \(\overline{f}^{1}([0, a])\) （其中 a 取遍一切 \(> 0\) 的实数集）为围域基本系的一致结构。

X 上的两个伪度量称为等价的，如果它们定义同一个一致结构。

注。1) 若 \((a_{n})\) 是任一 \(>0\) 且趋于 \(0\) 的数列，则诸 \(\mathrm{U}_{a_n}\) 构成由 \(f\) 定义的一致结构的一个围域基本系。

\begin{enumerate}
\def\labelenumi{\arabic{enumi})}
\setcounter{enumi}{1}
\tightlist
\item
  用伪度量 \(f\) 定义一致结构，就是取 \(f\) 关于“ \(o\) 在子空间 \([o, +\infty]\) （\(\overline{\mathbf{R}}\) 的子空间）中的邻域滤子”的原像作为围域基本系。注意这一程序与使我们得以定义拓扑群上的一致结构的程序十分相似（第三章，§ 3，第 1 号）。
\end{enumerate}

设 \(f\) 与 \(g\) 是 \(\mathbf{X}\) 上的两个伪度量。由定义 2 可知，由 \(f\) 定义的一致结构粗于由 \(g\) 定义的一致结构，当且仅当对每个 \(a > 0\) 存在 \(b > 0\) ，使得关系 \(g(x, y) \leqslant b\) 蕴含 \(f(x, y) \leqslant a\) 。\(f\) 与 \(g\) 是等价伪度量的必要充分条件是：对每个 \(a > 0\) 存在 \(b > 0\) ，使得 \(g(x, y) \leqslant b\) 蕴含 \(f(x, y) \leqslant a\) ，且 \(f(x, y) \leqslant b\) 蕴含 \(g(x, y) \leqslant a\) 。

特别地，若存在常数 \(k\) 使得 \(f \leqslant kg\) ，则由 \(f\) 定义的一致结构粗于由 \(g\) 定义的一致结构。

设 \(\varphi\) 是区间 \([o, +\infty]\) 到其自身的映射，满足下列条件：1) \(\varphi(o) = o\) ，且 \(\varphi\) 在 \(o\) 处连续；2) \(\varphi\) 在 \([o, +\infty]\) 中递增，且在 \(o\) 的一个邻域中严格递增；3) 对一切 \(u \geqslant o\) 与 \(v \geqslant o\) ，有 \(\varphi(u + v) \leqslant \varphi(u) + \varphi(v)\) 。那么，若 \(f\) 是集合 X 上的任一伪度量，则复合 \(g = \varphi \circ f\) 是与 \(f\) 等价的伪度量。

读者容易验证，例如可取 \(\varphi\) 为下列函数中的任何一个：

\[
\sqrt {u}, \quad \log (\mathrm{I} + u), \quad \frac {u}{\mathrm{I} + u}, \quad \inf (u, \mathrm{I}).
\]

最后两个例子表明，对任一给定的伪度量（无论有限与否），总存在与它等价的有界伪度量。

定义 3. 若 \((f_{t})_{t\in I}\) 是集合 \(X\) 上一族伪度量，则在 \(X\) 上由诸伪度量 \(f_{t}\) 定义的一致结构组成的集的上确界称为由族 \((f_{t})\) 定义的一致结构。

X 上的两个伪度量族称为等价的，如果它们在 X 上定义同一个一致结构。

由一致结构组成的集的上确界的定义（第二章，§ 2，第 5 号），由伪度量族 \((f_{t})_{t\in\mathbf{I}}\) 在 X 上定义的一致结构 U 的围域滤子，是由诸集 \(\overline{f}_{t}^{1}([o,a])\) （其中 t 取遍 I，a 取遍 \textgreater0 的实数集）所生成的滤子（第一章，§ 6，第 2 号）。换句话说，按下述方式可得到 U 的一个围域基本系：任意取有限多个指标 \(\iota_{1},\iota_{2},\ldots,\iota_{n}\) ，并对每个 \(\iota_{k}\) 取一个数 \(a_{k}>0\) ；然后考察由点对 \((x,y)\in\mathbf{X}\times\mathbf{X}\) （满足 \(f_{\iota_{k}}(x,y)\leqslant a_{k}\) ， \(i\leqslant k\leqslant n\) ）所组成的集 V；这些集 V（对 n、诸 \(\iota_{k}\) 与诸 \(a_{k}\) 的一切可能选取）构成 U 的一个围域基本系。此外，我们可以限于所有 \(a_{k}\) 都等于同一个数 a\textgreater0 的情形，因为由一切点对 \((x,y)\) （它们满足

\[
\sup _ {1 \leqslant k \leqslant n} \left(f _ {\iota_ {k}} (x, y)\right) \leqslant \inf _ {1 \leqslant k \leqslant n} a _ {k}
\]

）所组成的围域显然包含在 V 中。

对 I 的每个有限子集 H，用 \(g_{H}\) 表示族 \((f_{t})_{t\in\mathbb{H}}\) 的上包络。当 H 取遍 I 的所有有限子集、a 取遍 \textgreater0 的实数集时，诸集 \(g_{\mathrm{H}}^{-1}([o,a])\) 构成一致结构 U 的一个围域基本系。而诸 \(g_{H}\) 是 X 上的伪度量（第 1 号），且按定义，族 \((g_{\mathrm{H}})\) 中有限多个函数的上包络仍属于该族；我们把这一性质表述为：伪度量族 \((g_{\mathrm{H}})\) 是饱和的。于是伪度量族 \((g_{\mathrm{H}})\) 与族 \((f_{t})\) 等价，并称为由 \((f_{t})\) 饱和化而得的伪度量族。由刚才所述可知，我们总可以限于考虑由伪度量饱和族定义的一致结构。

在 \(\mathbf{I}\) 为有限集的特殊情形，这一论证表明，由伪度量族 \((f_{t})_{t\in \mathbf{I}}\) 定义的一致结构也可由单个伪度量 \(g = \sup_{t\in \mathbf{I}}f_t\) 定义。

设 \(\mathcal{U},\mathcal{U}'\) 是 \(\mathbf{X}\) 上的两个一致结构，分别由两个饱和族 \((f_{t})_{t\in \mathbf{I}},(g_{x})_{x\in \mathbf{K}}\) 定义。则 \(\mathcal{U}\) 粗于 \(\mathcal{U}'\) ，当且仅当对每个指标 \(t\in I\) 与每个实数 \(a > 0\) ，存在指标 \(x\in K\) 与数 \(b > 0\) ，使得关系 \(g_{x}(x,y)\leqslant b\) 蕴含 \(f_{t}(x,y)\leqslant a\) 。

由伪度量族定义的一致结构的例子。设 \((f_{t})_{t\in I}\) 是定义在集合 \(X\) 上的任意一族（有限）实值函数。设 \(\mathcal{U}\) 是 \(X\) 上使诸 \(f_{t}\) 都一致连续的最粗的一致结构（第二章，§ 2，第 3 号）。于是由 \(\mathcal{U}\) 的围域的定义（同前）可知，\(\mathcal{U}\) 就是 \(X\) 上由伪度量

\[
g _ {i} (x, y) = | f _ {i} (x) - f _ {i} (y) |.
\]

